\section{Structural congruence}

Processes are identified up to $\alpha$-conversion of restricted channels and of
bound variables. The structural congruence $\equiv$ is the least equivalence
relation containing the rules of \cref{tab:cgr} and closed under all
contexts: conditionals, parallel composition, restriction, recursion,
prefixes and sums. Congruent processes have the same tags.

\begin{table}[h]
\centering
\small
$\begin{array}{l@{\quad}rcl@{\quad}l}
\textsc{ParNil} & P \parallel \mathbf{0} &\equiv& P \\
\textsc{ParComm} & P \parallel Q &\equiv& Q \parallel P \\
\textsc{ParAssoc} & P \parallel (Q \parallel R) &\equiv& (P \parallel Q) \parallel R \\[2pt]
\textsc{SumNil} & K + \mathbf{0}_{\tilde e} &\equiv& K \\
\textsc{SumComm} & K + K' &\equiv& K' + K \\
\textsc{SumAssoc} & K + (K' + K'') &\equiv& (K + K') + K'' \\[2pt]
\textsc{IteT} & \ite{e}{P}{Q} &\equiv& P & e \Downarrow \mathit{tt} \\
\textsc{IteF} & \ite{e}{P}{Q} &\equiv& Q & e \Downarrow \mathit{ff} \\
\textsc{ValExpr} & P &\equiv& P[v/e] & e \Downarrow v \\[2pt]
\textsc{ResNil} & \mathbf{0}_{\tilde e} \setminus c &\equiv& \mathbf{0}_{\tilde e} \\
\textsc{ResOrd} & (P \setminus c) \setminus d &\equiv& (P \setminus d) \setminus c \\
\textsc{ResPar} & (P \parallel Q) \setminus c &\equiv& P \parallel (Q \setminus c) & c \notin \fc(P)
\end{array}$
\caption{Structural congruence. In \textsc{ValExpr}, $e$ is an expression of
an output, of a conditional, of a deadlock $\mathbf{0}_{\tilde e}$ or of a
quantum prefix, and $P[v/e]$ replaces it by its value. $\fc(P)$ is the set of
free channels of $P$.}
\label{tab:cgr}
\end{table}

\begin{lemma}[Congruence is a bisimulation]\label{lem:cgr-bisim}
If $P \equiv Q$ and $P \step{\lambda}{s} P'$ in the program LTS, then
$Q \step{\lambda}{s} Q'$ for some $Q' \equiv P'$.
\end{lemma}

The choice is preserved because the choice of a synchronisation is an
unordered pair: \textsc{ParComm} exchanges the two sides of the
synchronisation, not its choice.

\begin{remark}\label{rem:debruijn}
Scope extrusion \textsc{ResPar} needs the $\alpha$-conversion of the restricted
channel: $(R \setminus d) \parallel Q$ with $d$ free in $Q$ must first rename
$d$. As in VACCS, the mechanisation uses de Bruijn indices for restricted
channels and for the variables bound by inputs and measurements:
\textsc{ResPar} reads $(\nu P) \parallel Q \equiv \nu (P \parallel Q^{+})$,
where $Q^{+}$ shifts the indices of $Q$, \textsc{ResOrd} exchanges the two
innermost indices, and the substitution $P[Q/X]$ shifts $Q$ under each binder,
so that $\equiv$ is closed under recursion.
\end{remark}
